
% Chapter 4 File

\chapter{Reverse Transcriptases: Origin, Boundaries and Classification}
\label{chapter3}
\thispagestyle{empty}

% https://claude.ai/chat/1baca69a-335b-4adc-a4db-be125238b3ba

% c41c63d1-b030-4f01-8e91-b3a94f2dbfd9 (17 references) with 9745bb59-0cbf-48c2-990d-ff737eb0d487 

The framework set out in Section~\ref{sec:current-classification} is the one within which
retrons are currently described, and it is a considerable achievement: it identified a third
component, established its association with the RT under a statistic designed to be blind to
function and genomic context, and did so with controls capable of failing. It is also a
framework built on 1{,}928 sequences, in which five of the eleven RT clades contain no
experimentally characterised member and three yield no recoverable \textit{msr}--\textit{msd}
structure. For those clades the designation \emph{retron} rests on the RT sequence alone ---
on the presence of regions X and Y --- and therefore on the assumption that the features
originally used to define retron RTs are the features that in fact delimit them. Can we
confidently say those are retrons?

The difficulty runs deeper than taxonomic coverage, because most of the quantities on which
the classification depends are inherited rather than measured. The RT0--RT7 partition that
defines the alignable core was fixed by earlier work and carries no uncertainty estimate;
clade delimitation followed node support rather than an explicit algorithmic rule; the
thresholds of the association pipeline were chosen once and never varied; and enough
parameters go unstated that the analysis cannot be reproduced as run. None of this makes the
conclusions wrong. It makes them untested, and it means that a result and the convention that
produced it are at present indistinguishable.

% This chapter separates the two. Section~\ref{sec:rt-partition} re-examines the RT0--RT7
% partition, asking what sequence and structural evidence actually supports boundaries that are
% otherwise transmitted by citation. [Section~X] revisits clade delimitation under an explicit
% criterion, and [Section~Y] asks what the association analysis yields at present-day scale and
% under varied thresholds. The aim is not to discard the classification but to determine which
% of its elements survive once the decisions behind them are stated, varied, and given
% confidence intervals.


% \section{Provenance and physical basis of the RT0–RT7 partition}


% =====================================================================
%  Thesis section: the RT0-RT7 partition
%  Drop-in .tex. Requires in the preamble:
%    \usepackage{booktabs, tabularx, array, ragged2e, amssymb}
%    \usepackage{threeparttable}      % table notes tied to the table
%    \usepackage{rotating}            % sidewaystable (landscape)
%    \usepackage{longtable}           % Table 2 may break across pages
%    \usepackage[table]{xcolor}       % only for \rowcolor, optional
%    \usepackage{tikz}                % only for the optional provenance chain
%    \usetikzlibrary{positioning, arrows.meta}
%
%  Bibliography keys used here are defined in rt_partition_refs.bib
% =====================================================================

% ---- local helpers (put in the preamble if you prefer) ---------------
\newcolumntype{L}[1]{>{\RaggedRight\arraybackslash}p{#1}}
\newlength{\rtwd}   % usable width of a 4-column table with @{} on both ends
\setlength{\rtwd}{\dimexpr\textwidth-6\tabcolsep\relax}
\newcommand{\derived}{\textsc{derived}}          % boundaries determined in that work
\newcommand{\inherit}{\textsc{inherited}}        % taken over, with a source cited
\newcommand{\noprov}{\textsc{inherited}$^{\ast}$}% taken over, no source cited
\newcommand{\motif}[1]{\texttt{#1}}
% ---------------------------------------------------------------------

\section{The RT0--RT7 partition: inherited boundaries, sequence evidence
         and structural correlates}
\label{sec:rt-partition}

% ---------------------------------------------------------------------
\subsection{Origin and propagation of the partition}
\label{sec:rt-partition-origin}

The anatomy of the reverse transcriptase (RT) domain in prokaryotic retroelements rests
on a descriptive framework that has been remarkably stable since it was first proposed.
Xiong and Eickbush~\cite{xiong1990} aligned 82 retroelement RTs and identified seven
peptide regions---domains~1--7, comprising 178 amino acids common to all elements---and
it is this partition that almost every subsequent treatment inherits. Their scheme does
\emph{not} include RT0; that upstream element was described later in the non-LTR
retroelement literature~\cite{malik1999,zimmerly2001} and was integrated into the group~II
intron framework by Blocker \textit{et al.}~\cite{blocker2005}, who combined limited
proteolysis, Edman degradation and mass spectrometry with secondary-structure prediction
and homology modelling to produce the reference anatomy of the \textit{Lactococcus lactis}
Ll.LtrB maturase LtrA. That work placed RT1--RT7 within the fingers and palm subdomains of
the canonical polymerase fold, assigned RT0 to an extended fingers region, located
domain~X downstream of RT7 in the position corresponding to the thumb and part of the
connection domain of HIV-1 RT, and described the lineage-specific insertions 2a, 3a, 4a
and 7a relative to HIV-1 RT. Of these, 2a is not original to that work either: it
entered the framework alongside RT0 and from the same
sources~\cite{malik1999,zimmerly2001}, and both are shared by only a subset of RT
classes rather than being general features of the fold~\cite{simon2008}.

Table~\ref{tab:rt-sources} summarises what each source contributes and, in its final
column, records whether the element boundaries were determined in that work or adopted
from earlier literature. That column is an assessment made here by reading the published
methods; it is not a classification any of these authors apply to their own work. On this
reading the pattern is consistent: the boundaries were delimited by inspection of a
curated alignment~\cite{xiong1990}, no algorithmic criterion was published with them, and
the retron-specific lineage on which this thesis builds carries the span forward without
re-deriving it. The sole methodological statement in Toro and
Nisa-Mart{\'i}nez~\cite{toro2014} is that ``the amino-acid region corresponding to RT
domains 0 to 7'' was extracted---no criterion, and no source cited for the
span.\footnote{Neither Xiong and Eickbush~\cite{xiong1990} nor Blocker
\textit{et al.}~\cite{blocker2005} appears in its reference list, although
Mohr \textit{et al.}~\cite{mohr1993} is cited for domain~X. The absence is therefore
specifically of a source for the RT0--RT7 span, not of the Lambowitz-lab literature
generally. Stated as an absence from the published main text, since a definition could
reside in supplementary material.} Toro \textit{et al.}~\cite{toro2019} and
Rodr{\'i}guez Mestre \textit{et al.}~\cite{mestre2020} then use ``RT0--7'' as a given.
Both name individual elements when locating the retron-specific regions---region~X
between RT2 and RT3, region~Y within RT7---and neither refers to RT1 anywhere in its
published main text. The one exception to this pattern of inheritance is the myRT
reference package~\cite{sharifi2022}, whose profile models are built on the Pfam RVT\_1
domain rather than on the RT0--RT7 span; a coordinate system independent of the
partition is therefore already in use, though it has not been applied to the boundary
question.

A related inconsistency affects any attempt to reproduce these datasets. The
200-amino-acid minimum applied throughout this literature does not refer to a consistent
object: \cite{toro2014} discards proteins shorter than 200~aa as probable fragments,
applying the threshold to the \emph{full protein}, whereas~\cite{toro2019} selects
``an RT domain (RT0-7) of at least 200 amino acids'', applying it to the \emph{extracted
span}, a convention the retron datasets~\cite{mestre2020} inherit. The two thresholds are
not interchangeable---a domain of 200~aa can sit within a protein several times that
length---so the same nominal filter admits different sequence sets depending on which
paper is being followed. Either convention is defensible; conflating them is not. The
consequences are not confined to reproducibility: as Section~\ref{sec:rt-partition-elements}
shows, the span to which the threshold is applied opens with an element that has no
published alignable extent in retrons.

% =====================================================================
%  TABLE 1 -- source / evidence / provenance
%  Landscape. If your thesis is one-sided and you dislike sidewaystable,
%  swap \begin{sidewaystable} -> \begin{table} and use \footnotesize.
% =====================================================================
\begin{sidewaystable}
\centering
\begin{threeparttable}
\caption[Primary sources for the RT0--RT7 partition]{%
  \textbf{Primary sources for the RT0--RT7 partition, by class of evidence.}
  The final column distinguishes work in which element boundaries were determined from
  work in which they were adopted, and---among the latter---whether a source for the span
  is cited. That assignment is an assessment made here from the published methods, not a
  characterisation offered by the authors cited.}
\label{tab:rt-sources}
\footnotesize
\begin{tabular}{@{}L{2.9cm} L{3.1cm} L{2.2cm} L{6.4cm} L{2.6cm}@{}}
\toprule
\textbf{Source} & \textbf{Material / scale} & \textbf{Evidence} &
\textbf{Contribution to the partition} & \textbf{Boundaries} \\
\midrule

Xiong \& Eickbush~\cite{xiong1990}
  & 82 retroelement RTs
  & Alignment
  & Origin of the partition: domains 1--7, 178 conserved aa. RT0 not defined
  & \derived\tnote{a} \\[2pt]

Malik \textit{et al.}~\cite{malik1999}; Zimmerly \textit{et al.}~\cite{zimmerly2001}
  & non-LTR retroelements
  & Alignment
  & RT0 described upstream of domain 1
  & \derived \\[2pt]

Blocker \textit{et al.}~\cite{blocker2005}
  & LtrA (Ll.LtrB maturase)
  & Proteolysis, Edman, MS; homology model
% NEW
    & Reference anatomy. RT1--RT7 in fingers and palm; RT0 within a 10\,kDa
      N-terminal proteolytic fragment (M1--R85); RT0--RT7 within M1--R364;
      domain~X = thumb and part of the connection, downstream of RT7 and
      \emph{outside} the RT0--RT7 span; insertions 2a, 3a, 4a, 7a
    & \inherit~\cite{xiong1990,malik1999,zimmerly2001} \\[2pt]

Simon \& Zimmerly~\cite{simon2008}
  & 1{,}021 RTs; 20 groups
  & Alignment, comparative
  & First quantification of transferability: alignable characters per class in domains
    1--7 (group II 177, retron 157, DGR 179); 59-character core in domains 3--5
    alignable across all bacterial RTs; domains 1, 2 and 7 unalignable for many
    bacterial RTs
  & \inherit~\cite{xiong1990}\tnote{b,d} \\[2pt]

Toro \& Nisa-Mart{\'i}nez~\cite{toro2014}
  & 742 bacterial RTs
  & Phylogeny, motif survey
  & Reference bacterial RT phylogeny; class-specific motifs (retron \motif{NAXXH},
    \motif{VTG}; DGR \motif{Y(V/M)DD}; G2L4 \motif{YIDD})
  & \noprov \\[2pt]

Toro \textit{et al.}~\cite{toro2019}
  & 198{,}760 $\rightarrow$ 9{,}141 representatives
  & Phylogeny
  & The dataset the retron field inherits; ``RT0-7'' used as a given
  & \noprov \\[2pt]

Rodr{\'i}guez Mestre \textit{et al.}~\cite{mestre2020}
  & 1{,}912 retron / retron-like $+$ 16 validated
  & Phylogeny
  & Retron-specific reference; 11 clades; applies region~X (\motif{NAXXH},
    between RT2 and RT3) and region~Y (\motif{VTG}, within RT7)
  & \noprov\tnote{c} \\[2pt]

Yang \textit{et al.}~\cite{yang2025}
  & G2L4 RT, apo and DNA-bound (PDB 9D5X, 9D4S)
  & X-ray crystallography
  & Crystallographic resolution of the elements. RT0 is a loop within an N-terminal
    extension; NTE/RT0, RT2a and RT3a shared with RdRPs and lost from retroviral RTs;
    RT1 disordered in the apoenzyme
  & \textsc{structural} \\[2pt]

Arkhipova~\cite{arkhipova2023}; Gladyshev \& Arkhipova~\cite{gladyshev2011}
  & RT superfamily
  & Review; comparative
  & Places prokaryotic lineages in the wider fold; AbiP2 and AbiK lack the RT thumb,
    replaced by HEAT repeats
  & --- \\[2pt]

% NEW
Sharifi \& Ye~\cite{sharifi2022}
  & 1{,}844 seqs, 38 classes (FastTree placement package; the RVT-All.hmm
    release comprises 45 models for 41 classes from 1{,}988 sequences)
  & Profile HMM, phylogeny
  & myRT classifier and curated reference package used throughout this work.
    Models are built on the Pfam RVT\_1 domain (PF00078), \emph{not} on the
    RT0--RT7 span
  & \textsc{independent}\tnote{e} \\[2pt]

Khan \textit{et al.}~\cite{khan2025}
  & $>$100 previously untested retrons
  & Experiment
  & Functional readout (RT-DNA production, editing) against which sequence features are
    tested
  & --- \\

\bottomrule
\end{tabular}
\begin{tablenotes}[flushleft]\footnotesize
\item[$\ast$] Adopted without a cited source for the span, as far as the published main
  text shows.
\item[a] Delimited by inspection of the alignment; no algorithmic criterion was published,
  and no uncertainty is attached to the boundaries.
\item[b] Per-class character counts are, by the authors' own footnote, partly estimated by
  comparison with group II introns where a group lacked strong similarity in one or more
  domains.
\item[c] Regions X and Y are attributed to earlier retron work~\cite{simon2019} rather than
  derived. They were already in use as retron classification hallmarks in
  2008~\cite{simon2008}, citing Lampson \textit{et al.}~\cite{lampson2005}.
\item[d] The alignment presented spans domains 0--7, but every quantitative measure is
  taken over domains 1--7 and the cross-class core is located in domains 3--5. No
  transferability figure is therefore published for RT0 or 2a by the only study to
  quantify transferability.
\item[e] Does not adopt the RT0--RT7 partition: profile models are built on the Pfam
  RVT\_1 domain (PF00078). An existing precedent for a coordinate system that does not
  presuppose the partition.
\end{tablenotes}
\end{threeparttable}
\end{sidewaystable}

% ---------------------------------------------------------------------
% \subsection{Element-by-element: what is supported by sequence, and what by structure}
% \label{sec:rt-partition-elements}

% Table~\ref{tab:rt-elements} reorganises the same literature by element rather than by
% source, which is the view that matters for the boundary estimates developed in this
% chapter. Read down the two central columns, the partition separates into three groups. The
% catalytic core (RT3--RT5) is supported by both routes and carries the only universally
% alignable characters. The insertions (RT2a, RT3a) and RT0 are structurally well defined but
% poorly conserved in sequence. RT1 and RT7, the outer elements, are weak by both routes---and
% in the case of RT1 the crystallographic evidence supplies a mechanism.

% Yang \textit{et al.}~\cite{yang2025} report that RT1, which ordinarily forms two
% antiparallel $\beta$-strands including a fingertips loop over the active site, was
% disordered and not visible in the G2L4 apoenzyme, having been displaced by an RT3a plug
% occupying the active site; on DNA binding, extrusion of that plug allowed RT1 to re-form the
% expected strands. RT1 is therefore not a static block but a conformationally dynamic element
% whose order depends on the catalytic state of the enzyme. Yang \textit{et al.} also report the limit of this mechanism: in group~II intron
% RT apoenzymes obtained with constructs lacking a thumb (PDB 5HHJ, 5HHL) the
% shorter RT3a, which lacks a 9-amino-acid insertion present in G2L4, does not
% displace the RT1 strands, which remain ordered over the active site. They
% describe the insertion as a unique adaptation of G2L4. The displacement
% mechanism is therefore lineage-specific and does not by itself account for the
% unalignability of domain~1 across bacterial RTs reported
% earlier~\cite{simon2008}. Two readings remain open: that RT1 is conformationally
% labile in a way only sometimes captured, or that the sequence and structural
% weaknesses have separate causes. Distinguishing them requires tabulating
% unmodelled RT1 residues across deposited prokaryotic RT structures, which is
% done in Section~\ref{sec:results-rt1}. Neither reading is assumed here.

% ---------------------------------------------------------------------
\subsection{Element-by-element: what is supported by sequence, and what by structure}
\label{sec:rt-partition-elements}

Table~\ref{tab:rt-elements} reorganises the same literature by element rather than by
source, which is the view that matters for the boundary estimates developed in this
chapter; Table~\ref{tab:rt-occupancy} records, for the one published alignment in which
the partition is applied across classes, which elements could actually be aligned in
each. Read together, the partition separates into three groups rather than behaving as a
single object. The catalytic core, RT3--RT5, is supported by both routes: it carries the
59 characters alignable across all bacterial RTs in~\cite{simon2008}, and RT3 and RT5 are
the only two elements aligned in every bacterial representative of that study's Figure~4.
The insertions 2a and 3a are structurally well defined but lineage-restricted, and are
not elements of the RT0--RT7 span in the sense the other rows are. RT0, RT1 and RT7 are
weak by both routes, and each fails in a different way.

RT0 is the element the retron literature reports in and cannot align. It is absent from
the original partition~\cite{xiong1990}, was added later from the non-LTR
literature~\cite{malik1999,zimmerly2001}, and in LtrA is delimited only as the content of
a 10\,kDa proteolytic fragment~\cite{blocker2005}. It is excluded from every
transferability figure published by the study that quantified transferability, whose
counts are taken over domains 1--7 although its alignment is presented over
0--7~\cite{simon2008}. And in that alignment it is not alignable with the group~II
reference in any of the three retron or three DGR representatives shown. This is a
statement about the coordinate system and not about the proteins: the authors are
explicit that sufficient residues are present to correspond to the elements that cannot
be aligned. But it means that the span on which the retron datasets are
filtered~\cite{toro2019} and the retron phylogeny is inferred~\cite{mestre2020} opens
with an element that has no published alignable extent in retrons, and that the
$\ge$200-amino-acid threshold applied to that span therefore measures something whose
starting point is undefined for a quarter of the dataset.

RT1 and RT7 fail at the outer edges of the 1--7 span. Both were reported as unalignable
for many bacterial RTs~\cite{simon2008}; in the occupancy reading, RT1 is not aligned in
AbiA and only partially in AbiK, while RT7 is not alignable in any AbiK, AbiA or Abi-P2
representative and is partial or absent in the retrons. For RT1 the crystallographic
literature supplies a candidate mechanism and, in the same source, its limit. Yang
\textit{et al.}~\cite{yang2025} report that RT1, which ordinarily forms two antiparallel
$\beta$-strands including a fingertips loop over the active site, was disordered and not
visible in the G2L4 apoenzyme, having been displaced by an RT3a plug occupying the active
site; on DNA binding, extrusion of that plug allowed RT1 to re-form the expected strands.
They also report that in group~II intron RT apoenzymes obtained with constructs lacking a
thumb (PDB 5HHJ, 5HHL) the shorter RT3a, which lacks a nine-residue insertion present in
G2L4, does not displace the RT1 strands, which remain ordered over the active-site
pocket; the insertion is described there as a unique adaptation of G2L4. The displacement
mechanism is therefore lineage-specific as currently documented, and does not by itself
account for the unalignability of domain~1 across bacterial RTs reported
earlier~\cite{simon2008}. Two readings remain open: that RT1 is conformationally labile
in a way only sometimes captured crystallographically, of which G2L4 is the clearest
instance; or that the sequence and structural weaknesses of RT1 have separate causes and
their coincidence is not evidence of one. Distinguishing them requires tabulating which
RT1 residues are unmodelled across deposited prokaryotic RT structures rather than
reasoning from one pair, which is done in Section~\ref{sec:results-rt1}. Neither reading
is assumed in deriving the boundaries below.

% An element ordered only in the
% active state would be expected to align poorly in sequence and to be inconsistently resolved
% across structures. That is the pattern this work reports for RT1 from sequence and from
% deposited structures (Section~\ref{sec:results-rt1}), and it is consistent with the
% unalignability of domains 1, 2 and 7 reported earlier~\cite{simon2008}. This is stated as a hypothesis consistent with the evidence rather than
% as a demonstrated explanation; it is not tested in this work.

% % =====================================================================
%  TABLE 2 -- element x evidence matrix
% =====================================================================
\begingroup
\footnotesize
\setlength{\LTcapwidth}{\textwidth}
\begin{longtable}{@{}L{0.13\rtwd} L{0.18\rtwd} L{0.35\rtwd} L{0.34\rtwd}@{}}
\caption[Sequence and structural evidence per RT element]{%
  \textbf{Sequence and structural evidence, element by element.}
  Rows follow the conventional order of the RT0--RT7 span; domain X and the retron regions
  X and Y are appended because they are reported in the same coordinate system while lying
  partly outside it.}
\label{tab:rt-elements}\\
\toprule
\textbf{Element} & \textbf{Fold location}$^{a}$ & \textbf{Sequence evidence} &
\textbf{Structural evidence} \\
\midrule
\endfirsthead

\multicolumn{4}{@{}l}{\emph{Table \thetable\ (continued)}}\\
\toprule
\textbf{Element} & \textbf{Fold location} & \textbf{Sequence evidence} &
\textbf{Structural evidence} \\
\midrule
\endhead

\midrule
\multicolumn{4}{r@{}}{\emph{continued on the next page}}\\
\endfoot

\bottomrule
\addlinespace[2pt]
\multicolumn{4}{@{}p{\textwidth}@{}}{\footnotesize
  $^{a}$~Blocker \textit{et al.}~\cite{blocker2005} assign RT1--RT7 collectively to the
  fingers and palm subdomains; block-level assignments are approximate except for RT1,
  placed over the active site~\cite{yang2025}, and RT5, which carries the catalytic
  aspartates. No source consulted states element boundaries residue-by-residue against
  the subdomains, and the colouring in~\cite{yang2025}, Fig.~1A does not everywhere
  agree with the ordering assumed here; the assignments in this column should be read as
  approximate throughout.\newline
  $^{b}$~Note the nomenclature collision: \emph{domain}~X denotes the thumb and part of
  the connection subdomain~\cite{blocker2005}, whereas \emph{region}~X denotes the retron
  \motif{NAXXH} element~\cite{toro2014,mestre2020}. The two are unrelated and lie at
  opposite ends of the molecule. This thesis uses ``thumb'' for the former throughout,
  noting that the correspondence is to the thumb and part of the connection rather than
  to the thumb alone.\newline
  $^{c}$~Percentages in the regions X and Y rows are reported
  in~\cite{toro2014} over 102 retron/retron-like sequences at ${\le}85\%$ identity, with
  no matching rule stated for the motifs.}\\
\endlastfoot

RT0 / NTE & Extended fingers
  & Absent from~\cite{xiong1990}; described in non-LTR RTs~\cite{malik1999,zimmerly2001};
    within a 10\,kDa N-terminal proteolytic fragment of LtrA (M1--R85), conserved
    Ala39~\cite{blocker2005}; \emph{not alignable with the group~II reference in
    any of the three retron or three DGR representatives}~\cite{simon2008};
    excluded from all published transferability counts~\cite{simon2008}
  & Loop within an N-terminal extension~\cite{yang2025}. Helical character
    predicted in~\cite{blocker2005} \\[3pt]

RT1 & Fingertips, over the active site
  & Unalignable for many bacterial RTs~\cite{simon2008}; no reference to RT1 in
    \cite{toro2019,mestre2020} (inspection of the published text)
  & Two antiparallel $\beta$-strands when ordered; \emph{disordered} in the G2L4
    apoenzyme (9D5X), displaced by an RT3a plug; re-forms on DNA binding (9D4S).
    In group~II intron RT apoenzymes lacking the 9-aa RT3a insertion (5HHJ,
    5HHL) the RT1 strands remain ordered, and the insertion is described as a
    G2L4-specific adaptation~\cite{yang2025} \\[3pt]

RT2 & Fingers
  & Unalignable for many bacterial RTs~\cite{simon2008}
  & Resolved in both states~\cite{yang2025} \\[3pt]

RT2a & Insertion
  & LtrA \motif{SYGFRPX(K/R)S}, residues 130--138~\cite{blocker2005}
  & Shared with RdRPs; lost from retroviral RTs~\cite{yang2025} \\[3pt]

RT3 & Fingers--palm
  & Within the universal core: 59 alignable characters across domains
    3--5~\cite{simon2008}
  & Resolved in both states~\cite{yang2025} \\[3pt]

RT3a & Insertion
  & Lineage-specific; defined relative to HIV-1 RT~\cite{blocker2005}
  & Forms the plug occupying the active site in the apoenzyme; extruded on DNA
    binding~\cite{yang2025} \\[3pt]

RT4 & Palm & Universal core~\cite{simon2008}; insertion 4a~\cite{blocker2005}
  & Resolved in both states~\cite{yang2025} \\[3pt]

RT5 & Palm (catalytic)
    & \motif{Y/FXDD}; class-specific variants: \motif{YIDD} in G2L4~\cite{toro2014},
    \motif{Y(V/M)DD} in DGRs~\cite{toro2014}, conserved Arg at the variable position in
    Abi-P2~\cite{simon2008}, x nearly always hydrophobic
    & Catalytic \motif{YIDD} confirmed in G2L4~\cite{yang2025}; three invariant
    aspartates~\cite{arkhipova2023} \\[3pt]

RT6 & Palm & Conserved across classes~\cite{simon2008}
  & Resolved in both states~\cite{yang2025} \\[3pt]

RT7 & Palm, adjacent to the thumb\tnote{c}
  & Unalignable for many bacterial RTs~\cite{simon2008}; insertion 7a~\cite{blocker2005}
  & Not alignable in any AbiK or Abi-P2 representative, and partial or absent in
    the retron representatives~\cite{simon2008}. The Abi RT-like domain retains
    palm and fingers; it is the thumb that is replaced by a helical repeat
    domain~\cite{figiel2022}, which does not bear on RT7 \\[3pt]

\addlinespace
\multicolumn{4}{@{}l}{\emph{Reported in the same coordinate system but lying outside the
  RT0--RT7 span}}\\
\addlinespace

Domain X$^{b}$ & Thumb and part of the connection
  & Downstream of RT7; present within the 27\,kDa X-D-En proteolytic fragment
    (R365--K599, or S372--K599 from a second cleavage)~\cite{blocker2005}
  & Not resolved here \\[3pt]

Region X$^{b,c}$ & Between RT2 and RT3 &
  \motif{NAXXH}, identified in 59\% of 102 retron/retron-like sequences; His 100\%, Ala 87\%, Asn 62\% conserved~\cite{toro2014}; suggested to be required for initiation or catalysis~\cite{mestre2020} &
  Not resolved \\[3pt]
  
Region Y$^{b,c}$ & Within RT7 &
  \motif{VTG}, Gly most conserved; 47\% of the same set, with a further 26\% carrying an I/L variant at the first position~\cite{toro2014}; implicated in targeting the RT to its cognate \textit{msr}~\cite{inouye1999,mestre2020} &
  Not resolved \\

\end{longtable}
\endgroup


\newcommand{\algn}{\ensuremath{\bullet}}
\newcommand{\prtl}{\ensuremath{\circ}}
\newcommand{\nalg}{\textendash}

\begin{table}[htbp]
\centering
\begin{threeparttable}
\caption[Element occupancy by RT class]{%
  \textbf{Element occupancy by RT class, read from the published alignment.}
  Status of each element in the representative sequences of Figure~4
  of~\cite{simon2008}. \algn~aligned with the group~II reference,
  \prtl~partial or ambiguously aligned, \nalg~not alignable. Fractions give the
  number of representatives showing the majority state where a class is not uniform.}
\label{tab:rt-occupancy}
\footnotesize
\begin{tabular}{@{}l c *{9}{c}@{}}
\toprule
\textbf{RT class} & \textbf{$n$} &
\textbf{0} & \textbf{1} & \textbf{2} & \textbf{2a}\tnote{a} &
\textbf{3} & \textbf{4} & \textbf{5} & \textbf{6} & \textbf{7} \\
\midrule
\multicolumn{11}{@{}l}{\emph{Bacterial}}\\
Group II introns & 3 & \algn & \algn & \algn & \algn & \algn & \algn & \algn & \algn & \algn \\
Group II-like 1  & 4 & \algn & \algn & \algn & \algn & \algn & \algn & \algn & \algn & \algn \\
Group II-like 2  & 3 & \algn & \algn & \algn & \algn~(2/3) & \algn & \algn & \algn & \algn & \algn \\
Group II-like 3  & 2 & \algn & \algn & \algn & \nalg & \algn & \algn & \algn & \algn & \algn \\
Group II-like 4  & 2 & \prtl & \algn & \algn & \algn & \algn & \algn & \algn & \prtl & \algn \\
Group II-like 5  & 2 & \algn & \algn & \algn & \algn & \algn & \algn & \algn & \nalg & \algn \\
Retrons          & 3 & \nalg & \algn & \algn & \nalg & \algn & \algn & \algn & \algn & \prtl~(2/3) \\
DGRs             & 3 & \nalg & \algn & \algn & \algn & \algn & \algn & \algn & \algn & \algn \\
AbiK             & 3 & \algn & \prtl & \algn & \nalg & \algn & \algn & \algn & \algn & \nalg \\
AbiA             & 1 & \nalg & \nalg & \nalg & \nalg & \algn & \nalg & \algn & \algn & \nalg \\
Abi-P2           & 3 & \algn & \algn & \algn & \nalg & \algn & \algn & \algn & \algn & \nalg \\
\addlinespace
\multicolumn{11}{@{}l}{\emph{Eukaryotic and viral, shown for context}\tnote{b}}\\
Retroplasmids    & 3 & \algn & \algn & \algn & \nalg & \algn & \algn & \algn & \algn & \algn \\
non-LTRs         & 3 & \algn & \algn & \algn & \algn & \algn & \algn & \algn & \algn & \algn \\
LTRs             & 3 & \nalg & \prtl & \algn & \nalg & \algn & \algn & \algn & \algn & \algn \\
Hepadnaviruses   & 3 & \nalg & \algn & \prtl & \nalg & \algn & \algn & \algn & \algn & \algn \\
PLEs             & 3 & \nalg & \algn & \algn & \nalg & \algn & \algn & \algn & \algn & \algn \\
TERTs            & 3 & \nalg & \prtl & \algn & \nalg & \algn & \algn & \algn & \algn & \prtl \\
\bottomrule
\end{tabular}
\begin{tablenotes}[flushleft]\footnotesize
\item[$\ast$] The scoring is a reading of the published figure made here, not a
  classification offered by the authors. Following the figure legend, \nalg~indicates
  that the sequence could not be aligned with the group~II reference, \emph{not} that
  the residues are absent: the authors state that sufficient amino acids are present to
  correspond to the missing domains, with the possible exception of 2a. Three
  representatives per class are shown in the figure except where noted.
\item[a] 2a is an insertion shared by a subset of lineages rather than a canonical
  element of the RT0--RT7 span, and is shown as annotation only; it is not counted
  towards occupancy.
\item[b] The cross-class measures of~\cite{simon2008} are derived over bacterial RTs
  only. These rows are reproduced for comparison and contribute to no count reported
  here.
\end{tablenotes}
\end{threeparttable}
\end{table}


% ---------------------------------------------------------------------
\subsection{Consequence for the present work}
\label{sec:rt-partition-consequence}

Taken together, Tables~\ref{tab:rt-sources} and~\ref{tab:rt-elements} identify what is
measured and what is inherited. The boundaries carry no uncertainty estimate, the profile
models used to locate RT elements in new sequences are built from alignments constructed
around those same boundaries, and the coordinate system in which retron elements are
conventionally reported has never been assigned an explicit criterion. Distinguishing the
measured from the inherited therefore requires deriving element boundaries directly from
sequence data, with stated criteria and confidence intervals, in a coordinate system that
does not presuppose the partition being tested. That is the task of
Chapter~\ref{chap:boundaries}.

One further caution applies to landmark-based approaches. The dNTP-binding motif is
\motif{AQG} in G2L4 and \motif{PQG} in group II intron RT~\cite{yang2025}: the general form
is \motif{XQG}, and any search keyed on a literal \motif{PQG} is class-specific by
construction.


% =====================================================================
%  OPTIONAL: provenance chain as a figure. Delete if the prose suffices.
% =====================================================================
% \begin{figure}[htbp]
% \centering
% \footnotesize
% \begin{tikzpicture}[
%   node distance = 5mm and 0mm,
%   box/.style   = {draw, rounded corners=2pt, align=left, inner sep=4pt, text width=6.0cm},
%   note/.style  = {align=left, text width=5.2cm, font=\footnotesize\itshape},
%   arr/.style   = {-{Latex[length=2mm]}, thick}
% ]
% \node[box] (a) {Xiong \& Eickbush 1990 --- domains 1--7, 178 aa; no RT0};
% \node[box, below=of a] (b) {Malik 1999 / Zimmerly 2001 --- RT0 described};
% \node[box, below=of b] (c) {Blocker \textit{et al.} 2005 --- LtrA anatomy; RT0--RT7 $+$ thumb};
% \node[box, below=of c] (d) {Toro \& Nisa-Mart{\'i}nez 2014 --- 742-seq phylogeny};
% \node[box, below=of d] (e) {Toro \textit{et al.} 2019 --- 9{,}141 representatives};
% \node[box, below=of e] (f) {Rodr{\'i}guez Mestre \textit{et al.} 2020 --- retron clades; regions X, Y};

% \draw[arr] (a) -- (b); \draw[arr] (b) -- (c); \draw[arr] (c) -- (d);
% \draw[arr] (d) -- (e); \draw[arr] (e) -- (f);

% \node[note, right=4mm of a] {delimited by inspection};
% \node[note, right=4mm of c] {blocks drawn ``as defined by'' 1990, 2001};
% \node[note, right=4mm of d] {200 aa applied to the \textbf{protein}; no source cited for the span};
% \node[note, right=4mm of e] {200 aa applied to the \textbf{span}};
% \node[note, right=4mm of f] {inherits the span; regions X, Y attributed to Simon 2019};
% \end{tikzpicture}
% \caption[Provenance of the RT0--RT7 span]{%
%   \textbf{Provenance of the RT0--RT7 span.} Two threads run through this literature and
%   only one carries provenance: Blocker \textit{et al.} state the source of their blocks
%   explicitly, whereas the Toro--Mestre lineage, on which the retron field relies, cites no
%   source for the span. The annotations record what the published methods state, and the
%   contrast drawn between the two threads is an observation made here.}
% \label{fig:rt-provenance}
% \end{figure}